% stress-index-methodology.tex
% ----------------------------------------------------------------------
% Methodology note for the Macro Pulse Bangladesh Composite Stress Index
%
% Compile with:  pdflatex stress-index-methodology.tex
%                (run twice for table of contents)
% Output: stress-index-methodology.pdf
% Link to the PDF from the website footer and from the index chart
% footnote.

\documentclass[11pt, a4paper]{article}

\usepackage[margin=1in]{geometry}
\usepackage{amsmath, amssymb}
\usepackage[T1]{fontenc}
\usepackage{lmodern}
\usepackage[protrusion=true,expansion=false]{microtype}
\usepackage{booktabs}
\usepackage{array}
\usepackage{enumitem}
\usepackage{hyperref}
\usepackage[hang,flushmargin]{footmisc}
\usepackage{parskip}
\usepackage{titlesec}
\usepackage{xcolor}

\definecolor{linkcolor}{HTML}{0a6640}
\hypersetup{
  colorlinks = true,
  linkcolor  = linkcolor,
  urlcolor   = linkcolor,
  citecolor  = linkcolor
}

\titleformat{\section}{\normalfont\Large\bfseries}{\thesection}{1em}{}
\titleformat{\subsection}{\normalfont\large\bfseries}{\thesubsection}{1em}{}

\title{\textbf{Bangladesh Macro Pulse: \\
       A Composite Stress Index from News-Derived Signals} \\
       \vspace{0.4em}
       \large Methodology Note, Version 1.0}

\author{Syed Abul Basher \\
        \small Gold River Analytics, Dhaka \\
        \small \href{https://syedbasher.github.io/bangladesh-early-warning-signals/}%
                    {syedbasher.github.io/bangladesh-early-warning-signals}}

\date{\today}

\begin{document}

\maketitle
\vspace{-1em}
\hrule
\vspace{1em}

\begin{abstract}
\noindent
This note documents the construction of the \textit{Macro Pulse Composite
Stress Index} for the Bangladesh economy. The index condenses structured
economic signals --- derived from automated classification of
Bangladesh-focused news --- across sixteen channels into a single daily
$0$--$100$ reading of macroeconomic stress. We describe the signal
generation pipeline, the per-channel intensity construction, the
percentile-rank normalisation strategy chosen for a short training
window, and the weighting scheme. We close with an explicit discussion
of limitations and a roadmap for anchoring each qualitative channel to
high-frequency observable data.
\end{abstract}

\vspace{0.5em}
\tableofcontents
\vspace{1em}
\hrule

\section{Purpose}

The Macro Pulse Composite Stress Index is intended as a lightweight
\textit{early-warning summary statistic} for analysts, investors, and
policy observers tracking the Bangladesh economy. It is not a forecast
of any particular variable, nor a substitute for Bangladesh Bank's
official monetary and balance-of-payments releases. Its purpose is to
compress a large volume of qualitative economic reporting into a single
daily reading and to flag \textit{relative} escalation or moderation in
macro conditions across sixteen policy-relevant channels.

\section{Data}

\subsection{Source: the signal pipeline}

The index is computed from \texttt{data/signals.json}, the structured
output of the Macro Pulse pipeline. Each signal $i$ in that file
comprises, at minimum, the fields
\[
  \bigl(
    \text{date}_i,\;
    \text{channel}_i,\;
    \text{direction}_i \in \{\text{tightening}, \text{easing}\},\;
    \text{confidence}_i \in \{\text{low}, \text{medium}, \text{high}\}
  \bigr).
\]
These fields are generated by an LLM classifier (Claude Haiku) applied
to Bangladesh-focused RSS headlines collected twice daily. The
classifier is governed by a fixed prompt documented in
\texttt{interpret\_events.py}; outputs are validated against a
controlled vocabulary and cached.

\subsection{Channels}

Sixteen structural channels cover the main macroeconomic transmission
mechanisms for Bangladesh:
Foreign Exchange, Exports, Remittances, Food Prices, Credit, Banking,
Energy, Monetary Policy, Fiscal, Political Risk, Agriculture,
Logistics, Trade Policy, Labor Market, Geopolitical Risk, and
Capital Markets. Channel definitions and metadata (lead indicator,
time horizon, economic mechanism, expected effects) are fixed in
\texttt{signal\_definitions.py}.

\subsection{Sign convention}

Signal direction encodes whether the reported condition worsens a
macroeconomic constraint (\textit{tightening}; e.g.\ reserves falling,
export orders cancelled, load shedding intensifying, NPL rising) or
relaxes it (\textit{easing}; e.g.\ reserves replenished, rate cut
transmitted, food-price stabilisation). We map
\[
  d_i = \begin{cases}
    +1 & \text{if direction}_i = \text{tightening} \\
    -1 & \text{if direction}_i = \text{easing}
  \end{cases}
\]
so that the per-channel intensity $I_c(t)$ defined below has the
property that higher values indicate more macro stress.

\section{Construction}

\subsection{Per-channel intensity}

For channel $c$ on date $t$, the \textit{raw intensity} is
\begin{equation}
  I_c(t) \;=\; \sum_{\substack{i \in c \\ 0 \le t - t_i \le H}}
               w(\text{conf}_i) \cdot d_i \cdot
               \exp\!\left(-\ln 2 \cdot \frac{t - t_i}{\tau}\right),
  \label{eq:intensity}
\end{equation}
where $H = 60$ days is the trailing window, $\tau = 14$ days is the
half-life governing exponential decay, and confidence weights are
\[
  w(\text{low}) = 1, \quad w(\text{medium}) = 2, \quad w(\text{high}) = 3.
\]
The decay term $2^{-(t - t_i)/\tau}$ ensures that a signal's
contribution halves every $\tau$ days. Older signals fade
continuously rather than drop off sharply at a fixed cut-off.

\paragraph{Choice of $\tau$.} A 14-day half-life balances two concerns.
Too short ($\tau \le 7$) produces a twitchy index dominated by the most
recent news cycle; too long ($\tau \ge 30$) smooths the series into a
trend indicator that reacts slowly to escalation. Fourteen days is also
consistent with the typical ``lead time'' between news-reportable
events (e.g.\ an LC crunch in the press) and observable macroeconomic
flow (e.g.\ monthly remittance settlements).

\paragraph{Role of the outer window $H$.} The $60$-day window is
redundant given exponential decay (by day 60, a signal's weight is
$2^{-60/14} \approx 5\%$) but is retained for computational
predictability and to guarantee that events beyond two months cannot
contribute mechanically.

\subsection{Percentile-rank normalisation}

Raw intensities are not directly comparable across channels because
(i) channels differ in their natural news volume (Energy generates many
more headlines than Agriculture in a typical week) and (ii) the short
training sample precludes stable variance estimation. We therefore
map each raw intensity to its empirical percentile rank:
\begin{equation}
  S_c(t) \;=\; 100 \cdot \hat{F}_{c}\bigl(I_c(t)\bigr),
  \label{eq:percentile}
\end{equation}
where $\hat{F}_c(\cdot)$ is the empirical CDF of $\{I_c(\tau) : \tau \le
t,\; I_c(\tau) \neq 0\}$, estimated with the mid-rank convention for
ties. Days on which channel $c$ has zero activity are assigned
$S_c(t) = 50$ (treated as neutral, not as ``historically low stress'').

Percentile-rank has three properties we want for this application:
\begin{enumerate}[leftmargin=2em, topsep=0.2em, itemsep=0.2em]
  \item It is distribution-free and robust to the small sample size.
  \item It is automatically self-calibrating as history accumulates.
  \item It is bounded in $[0, 100]$ without clipping.
\end{enumerate}

\subsection{Composite index}

The composite stress index $S(t)$ is a weighted average of per-channel
scores:
\begin{equation}
  S(t) \;=\; \frac{\sum_{c} \omega_c \cdot S_c(t)}{\sum_{c} \omega_c},
  \qquad S(t) \in [0, 100].
  \label{eq:composite}
\end{equation}

Channel weights $\omega_c$ reflect structural importance to the
Bangladesh macroeconomy, informed by the external-sector orientation of
the economy and the historical frequency with which each channel has
acted as a binding constraint:

\begin{center}
\small
\begin{tabular}{lr @{\hspace{2em}} lr}
\toprule
\textbf{Channel} & $\omega_c$ & \textbf{Channel} & $\omega_c$ \\
\midrule
Foreign Exchange     & 15 & Fiscal             & 6 \\
Exports              & 12 & Political Risk     & 6 \\
Remittances          & 10 & Agriculture        & 4 \\
Food Prices          & 10 & Logistics          & 3 \\
Credit               &  8 & Trade Policy       & 1 \\
Banking              &  8 & Labor Market       & 1 \\
Energy               &  8 & Geopolitical Risk  & 1 \\
Monetary Policy      &  7 & Capital Markets    & 0 \\
\midrule
\multicolumn{4}{r}{\textit{Total weight} = 100} \\
\bottomrule
\end{tabular}
\end{center}

Capital Markets is retained in the signal taxonomy but receives zero
weight in the composite; it is tracked separately via a hard-data
anchor (DSEX and portfolio flow data). Trade Policy, Labor Market, and
Geopolitical Risk receive low weight because their signals are
frequently second-order or already captured by the FX, Energy, and
Exports channels.

\subsection{Stress zones}

For interpretive clarity the $[0, 100]$ composite is binned into three
zones:

\begin{center}
\begin{tabular}{lcc}
\toprule
\textbf{Zone} & \textbf{Range} & \textbf{Interpretation} \\
\midrule
Green  & $S(t) \le 35$   & Stress historically muted \\
Amber  & $35 < S(t) < 65$ & Stress near historical median \\
Red    & $S(t) \ge 65$   & Stress historically elevated \\
\bottomrule
\end{tabular}
\end{center}

These thresholds are heuristic and are expected to be revisited once
the index has accumulated at least twelve months of history and the
underlying signal distributions have stabilised.

\section{Publication schedule}

The pipeline runs twice daily on a scheduled GitHub Action
(\texttt{.github/workflows/collect.yml}, 00:00 and 12:00 UTC). Each run
recomputes \texttt{data/signals.json}, recomputes
\texttt{data/stress\_index.json}, and commits both artefacts. The
public website reads the JSON directly; no manual step intervenes
between data generation and publication.

\section{Limitations}

Transparency about limitations is an essential part of the product, not
a disclaimer appended to it.

\paragraph{1. Classification noise.} The underlying signal
classifications are produced by a large language model operating on
headline text only. Misclassifications occur, particularly around
(a) the direction of news containing conflicting cues (``crisis
averted'', ``growth despite headwinds''), (b) very short or elliptical
headlines, and (c) signals that should cross-tag across channels.
Prompt iteration and cache cleaning are ongoing.

\paragraph{2. Source selection.} The pipeline ingests a fixed set of
RSS feeds, weighted toward major Bangladesh English-language outlets
(\textit{The Daily Star}, \textit{Business Standard},
\textit{Financial Express}, \textit{Prothom Alo} English) and
international wires. Bangla-only reporting is under-represented. No
single source should be treated as authoritative; the index's validity
depends on the aggregate.

\paragraph{3. Headline truncation.} Classification operates on
headlines, not full articles. Magnitude is not extracted: a headline
stating ``reserves fall'' and one stating ``reserves fall to
twenty-year low'' currently receive identical tagging subject only to
confidence-level differences.

\paragraph{4. Short training sample.} As of the methodology version date
the pipeline has accumulated approximately seven weeks of continuous
history. The percentile-rank normalisation is therefore calibrated
against a very short distribution. Readings will shift as the baseline
widens; readers should interpret the absolute level with
corresponding caution during this initial period. Relative movement ---
direction and magnitude of change from week to week --- is considerably
more reliable than the absolute level.

\paragraph{5. Weighting is structural, not estimated.} Channel weights
$\omega_c$ are prior beliefs informed by the structure of the
Bangladesh economy, not estimates from a supervised calibration
exercise. A data-driven re-weighting exercise --- regressing the
index on realised outcomes such as quarterly reserves change,
DSEX drawdowns, or reported GDP surprises --- will become feasible
once twelve to eighteen months of history have accumulated.

\paragraph{6. Publication lag is zero.} Unlike official statistics the
index has no publication lag, but this also means it contains no
revision cycle. Readings for past days will not change; the
percentile-rank baseline grows but previous entries are fixed.

\section{Roadmap}

Planned improvements (in order of priority):
\begin{enumerate}[leftmargin=2em, topsep=0.2em, itemsep=0.3em]
  \item \textbf{Hard-data anchors.} Each qualitative channel to be
        paired with one observable high-frequency series --- weekly
        gross reserves (FX), DSEX and portfolio flows (Capital
        Markets), Tk/USD interbank rate (FX), TCB retail food prices
        (Food Prices), call-money rate (Banking), monthly EPB exports
        (Exports), monthly BB remittance inflow (Remittances).
  \item \textbf{Track record page.} A curated archive of past episodes
        with the state of the index in the two to six weeks
        preceding, including honest documentation of misses.
  \item \textbf{Classifier audit.} Periodic manual review of a random
        sample of classifications (target 200 per quarter) with a
        published accuracy metric.
  \item \textbf{Empirical re-weighting.} Once twelve months of history
        is available, channel weights are to be revised via a
        supervised procedure targeting observable stress outcomes.
  \item \textbf{Bangla-language ingestion.} Extend the collector to
        Bangla-only outlets (\textit{Prothom Alo} Bangla,
        \textit{Bonik Barta}, \textit{Daily Kaler Kantho}).
\end{enumerate}

\section{Reproducibility}

The full pipeline is open source. The repository at
\url{https://github.com/SyedBasher/bangladesh-early-warning-signals}
contains the classification code, channel definitions, computation
script, and this methodology note. The canonical implementation of
Equations~\eqref{eq:intensity}--\eqref{eq:composite} is in
\texttt{scripts/compute\_stress\_index.py}. The signal corpus itself,
from which the index is derived, is published at
\texttt{data/signals.json}.

\section*{Contact}

Correspondence on methodology, licensing, or institutional access
should be directed to Syed Abul Basher at Gold River Analytics.

\vspace{1em}
\hrule
\vspace{0.5em}
\begin{center}
\footnotesize
This methodology note will be versioned. Material revisions are
flagged in the header; minor edits (typography, phrasing) are
incorporated silently.
\end{center}

\end{document}
